\chapter{CSF (Cerebrospinal Fluid) Analysis}
\section*{What Is This Test?}
Cerebrospinal fluid (CSF) analysis is a laboratory evaluation of the fluid that surrounds the brain and spinal cord within the subarachnoid space. CSF is produced primarily by the choroid plexus in the lateral, third, and fourth ventricles at a rate of approximately 500 mL/day, with a total volume of 125--150 mL. CSF analysis is essential for the diagnosis of: bacterial meningitis, viral meningitis/encephalitis, fungal meningitis (cryptococcal), tuberculous meningitis, subarachnoid hemorrhage, multiple sclerosis, Guillain-Barre syndrome (GBS, albumino-cytologic dissociation --- elevated protein with normal WBC), leptomeningeal carcinomatosis, CNS lymphoma, and normal pressure hydrocephalus (NPH). The analysis includes: opening pressure (measured with a manometer before CSF collection), macroscopic appearance (clear, cloudy/turbid, xanthochromic/yellow, bloody), cell count with differential (RBCs and WBCs), glucose (compared with concurrent serum glucose, CSF/serum glucose ratio), protein, Gram stain and bacterial culture (aerobic and anaerobic), and specialized tests as indicated (viral PCR panel, VDRL for neurosyphilis, oligoclonal bands/IgG index for multiple sclerosis \cite{freedman2005recommended}, cytology for malignant cells, angiotensin-converting enzyme [ACE] for neurosarcoidosis, cryptococcal antigen, and 14-3-3 protein and RT-QuIC for Creutzfeldt-Jakob disease) \cite{seehusen2003cerebrospinal}.
\section*{Alternative Names}
CSF Analysis; Spinal Fluid Analysis; Lumbar Puncture (LP) Fluid; Cerebrospinal Fluid Examination; Spinal Tap Fluid
\section*{Principle}
CSF is collected by lumbar puncture (LP) at the L3-L4 or L4-L5 interspace (below the termination of the spinal cord/conus medullaris). Four sterile tubes are collected sequentially: Tube 1: chemistry/protein and glucose; Tube 2: microbiology (Gram stain, culture); Tube 3: cell count with differential; Tube 4: cytology or special studies. Cell count is performed manually on a hemocytometer. Differential is performed on a cytocentrifuged preparation stained with Wright-Giemsa. Glucose and protein are measured on automated clinical chemistry analyzers. The traumatic LP (bloody tap) presents the challenge of distinguishing blood from a true subarachnoid hemorrhage versus traumatic puncture. In a traumatic LP: the RBC count decreases from tube 1 to tube 4; xanthochromia is absent; the CSF WBC/RBC ratio approximates that of peripheral blood. In a true subarachnoid hemorrhage: the RBC count is constant across all tubes; xanthochromia (yellow-orange discoloration from oxyhemoglobin and bilirubin) develops 2--4 hours after hemorrhage and persists for 2--4 weeks; the WBC count may be elevated due to chemical meningitis from blood breakdown products.
\section*{History}
Heinrich Quincke first described lumbar puncture in 1891, initially as a therapeutic maneuver to lower intracranial pressure in hydrocephalus, and the technique rapidly became a diagnostic mainstay for meningitis and other neurologic diseases. By the early 1900s, CSF cell counts and protein measurement were being used to characterize meningitis, and the Wassermann test was applied to CSF in 1906 for the diagnosis of neurosyphilis. The antibiotic era of the 1930s--1940s made rapid Gram stain and culture of CSF critically important for early targeted therapy. In the 1980s, detection of oligoclonal bands and an elevated IgG index became standard laboratory support for the diagnosis of multiple sclerosis. Modern CSF analysis has been extended by multiplex PCR panels for infectious pathogens and by specialized biomarkers such as 14-3-3 protein and RT-QuIC for Creutzfeldt-Jakob disease, yet opening pressure measurement and sequential tube collection remain essentially unchanged since Quincke's original procedure.
\section*{Why This Test Is Ordered}
\begin{itemize}
  \item Suspected meningitis or encephalitis (fever, headache, neck stiffness, altered mental status, photophobia) to confirm infection and identify the organism \cite{tunkel2004practice}
  \item Suspected subarachnoid hemorrhage when noncontrast head CT is negative or equivocal \cite{perry2013clinical}
  \item Evaluation of multiple sclerosis: CSF oligoclonal bands and IgG index supplement MRI and clinical criteria
  \item Suspected Guillain-Barre syndrome: albumino-cytologic dissociation (elevated protein with normal cell count)
  \item Suspected leptomeningeal carcinomatosis or CNS lymphoma: CSF cytology and flow cytometry
  \item Evaluation of opportunistic infections in immunocompromised patients (cryptococcal meningitis, tuberculosis)
  \item Suspected Creutzfeldt-Jakob disease: CSF 14-3-3 protein and RT-QuIC
  \item Assessment of normal pressure hydrocephalus: opening pressure measurement and response to CSF drainage
\end{itemize}
\section*{Primary vs Secondary Test}
For suspected acute CNS infection (meningitis, encephalitis) and subarachnoid hemorrhage, CSF analysis is a primary, often urgent diagnostic test whose results guide empiric therapy. For multiple sclerosis, Guillain-Barre syndrome, and malignancy it is a secondary confirmatory test used together with MRI, neurophysiologic studies, and clinical criteria.
\section*{How This Test Is Performed}
CSF is collected by lumbar puncture (LP) performed at the bedside or fluoroscopically, in the lateral decubitus or seated position, at the L3-L4 or L4-L5 interspace. Opening pressure is measured with a manometer before fluid is withdrawn. Four sterile tubes are collected sequentially (tube 1: chemistry; tube 2: microbiology; tube 3: cell count with differential; tube 4: cytology or special studies), totaling approximately 10--15 mL. Cell counts are performed manually on a hemocytometer, differentials on a cytocentrifuged Wright-Giemsa stain, and glucose and protein on automated chemistry analyzers. Cell count, glucose, protein, and Gram stain results are typically available within 1--2 hours; cultures require 24--72 hours.
\section*{What Preparation Is Needed}
No fasting or specific dietary preparation is required. Before the procedure, the clinician reviews the coagulation status and medications and determines whether a head CT is needed to exclude a mass lesion or increased intracranial pressure. Informed consent is obtained, and the patient should empty the bladder and be positioned comfortably. The operator confirms the landmarks and uses full sterile technique.
\section*{Contraindications and Precautions}
\begin{itemize}
  \item Absolute contraindications: suspected increased intracranial pressure from a supratentorial mass (risk of herniation), skin infection at the puncture site, severe coagulopathy or thrombocytopenia (platelet count typically below 50,000/mcL), and ongoing anticoagulation
  \item Relative contraindications: spinal epidural abscess, severe spinal stenosis, prior lumbar surgery, and uncooperative patients
  \item Head CT before LP is recommended when papilledema, focal neurologic deficits, new-onset seizures, or immunocompromise suggest a mass lesion
  \item A traumatic tap can confound the cell count and protein; comparison of tube 1 and tube 4 RBC counts and xanthochromia help distinguish it from subarachnoid hemorrhage
  \item Empiric antibiotics should not be delayed while awaiting LP in suspected bacterial meningitis; cultures remain positive for hours after antibiotics
\end{itemize}
\section*{Risks}
Post-lumbar puncture headache occurs in up to 10--30\% of patients and is usually self-limited, lasting 1--3 days; use of a small-gauge or atraumatic needle reduces its frequency. Less common risks include local pain, bleeding or infection at the puncture site, transient paresthesias, and, rarely, nerve root injury. Brain herniation is a rare but catastrophic complication that is largely preventable by adherence to contraindications and pre-LP imaging.
\section*{What Happens After the Test}
The patient is asked to lie flat for 1--4 hours and to drink fluids to reduce the risk of post-LP headache. Preliminary results (cell count, glucose, protein, Gram stain) are usually available the same day and guide empiric therapy; culture results follow in 24--72 hours. Positive findings prompt directed treatment and additional testing, such as neuroimaging, blood cultures, or serial CSF studies.
\section*{Understanding Results}
\subsection*{Bacterial Meningitis}
Opening pressure elevated (20--50 cm H2O). Appearance: cloudy/turbid/purulent. WBC: markedly elevated (1,000--10,000+ cells/mcL), neutrophilic predominance (>80\%). Glucose: markedly decreased (<40 mg/dL, CSF/serum ratio <0.4). Protein: markedly elevated (100--500+ mg/dL). Gram stain: positive in 60--90\% (organism-dependent). Culture positive \cite{vandebeek2004clinical}. Common organisms: S. pneumoniae, N. meningitidis, H. influenzae (unvaccinated), L. monocytogenes (elderly, immunocompromised), S. agalactiae/Group B Streptococcus (neonates).
\subsection*{Viral Meningitis/Encephalitis}
Opening pressure normal or mildly elevated. Appearance: clear. WBC: mildly to moderately elevated (10--500 cells/mcL), lymphocytic predominance (after 24--48 hours; may be neutrophilic early). Glucose: normal (>60\% of serum). Protein: normal or mildly elevated (50--100 mg/dL). Gram stain and bacterial culture negative. PCR for enterovirus, HSV-1/2, VZV, EBV, CMV.
\subsection*{Fungal/Tuberculous Meningitis}
WBC: moderately elevated, lymphocytic predominance. Glucose: markedly decreased (similar to bacterial). Protein: markedly elevated. India ink and cryptococcal antigen for Cryptococcus. AFB smear and culture, PCR for M. tuberculosis.
\subsection*{Subarachnoid Hemorrhage}
Bloody CSF in all tubes (constant RBC count). Xanthochromia after 2--4 hours (spectrophotometric or visual detection). Elevated protein from blood.
\subsection*{Multiple Sclerosis}
Normal opening pressure, clear. Normal WBC (mild lymphocytic pleocytosis in 30\%, <50 cells). Normal glucose. Normal or mildly elevated protein. Oligoclonal bands (OCB) present in CSF but absent from paired serum (2+ unique bands in CSF = positive) OR IgG index >0.7.
\subsection*{Guillain-Barre Syndrome}
Albumino-cytologic dissociation: markedly elevated protein (100--1,000 mg/dL) with normal WBC count (<5 cells/mcL). This pattern is characteristic.
\section*{Normal Results}
Opening pressure: 6--20 cm H2O (obese patients up to 25 cm). Appearance: clear, colorless. WBC: 0--5/mcL (mononuclear cells, lymphocytes). RBC: 0--5/mcL. Glucose: 50--80 mg/dL (approximately 60--70\% of concurrent serum glucose, normally >45 mg/dL absolute). Protein: 15--45 mg/dL (lumbar CSF; higher in cisternal and ventricular CSF). Gram stain and culture: no organisms. IgG index: 0.3--0.7. Oligoclonal bands: negative (no unique bands in CSF compared to serum).
\section*{Follow-up Tests}
Meningitis/encephalitis multiplex PCR panel (BioFire FilmArray ME Panel detects 14 pathogens simultaneously), blood cultures, CT head before LP if papilledema, focal neurologic deficits, new-onset seizures, or immunocompromised (to exclude a mass lesion/risk of herniation), CSF VDRL (neurosyphilis), CSF ACE (neurosarcoidosis), CSF cytology and flow cytometry (malignancy), CSF 14-3-3 protein and RT-QuIC (CJD), and serum IgG/albumin for IgG index calculation.
\section*{References}
\bibliographystyle{unsrtnat}
\bibliography{references}

\clearpage
\section*{Regulatory Status and Authorization Date}
Regulatory status applies to the specific assay, instrument, manufacturer, and intended use rather than to the test name alone. No single approval date can be assigned to this test category. For a commercial assay, verify the current product in the relevant regulator's database: the U.S. Food and Drug Administration (FDA) clearance, approval, or authorization; India's Central Drugs Standard Control Organisation (CDSCO) licence or permission; the European Union In Vitro Diagnostic Regulation (EU IVDR) CE certificate issued through the applicable conformity-assessment route; or the United Kingdom Medicines and Healthcare products Regulatory Agency (MHRA) registration and UKCA/CE status. Laboratory-developed tests may not have an individual FDA, CDSCO, EU IVDR, or MHRA approval date.
\clearpage
